% ==============================================================================
% ПАРТИЯ B03-044: Мысли в слух: Часть 1 (блоки DISS-005-B2225 .. DISS-005-B2265)
% Точный перенос рукописи v0.5 без изменений
% ==============================================================================

% DISS-005-B2225
\chapter*{Мысли в слух}
\addcontentsline{toc}{chapter}{Мысли в слух}
\label{ch:notes_informal}

% DISS-005-B2226
Данный раздел сформирован для пометки и оформления осеняющих меня мыслей, которые потом можно переносить в основной раздел. Так как я делаю это всё с гиперссылками, то проблем возникнуть не должно.\par\vspace{0.3cm}

% DISS-005-B2227
\par\vspace{0.2cm}

% DISS-005-B2228
\section*{ИИ}

% DISS-005-B2229
\subsection*{Общие знания про ML}

% DISS-005-B2230
Добавить описание и анализ научных статей, которые я успел изучить на английском\par\vspace{0.3cm}

% DISS-005-B2231
Подумать о добавлении доклада по философии\par\vspace{0.3cm}

% DISS-005-B2232
\par\vspace{0.2cm}

% DISS-005-B2233
Смотрю серию лекций «Полная история GPT\commentnote{DISS-005-C0140}{Алексей Евдаков}{Конкретно сейчас видео ютуб\par\url{https://www.youtube.com/watch?v=J5nGTnv3IN8}}| Котенков Игорь\commentnote{DISS-005-C0141}{Алексей Евдаков}{\url{https://www.youtube.com/@stalkermustang}}\commentnote{DISS-005-C0142}{Алексей Евдаков}{Так же у него есть телеграмм канал, с которого черпаю очень много информации, которые описывал в докладе по философии\par\url{https://t.me/seeallochnaya}}», появилось очень много общих полезных знаний. Начну их расписывать, пока не забыл:\par\vspace{0.3cm}

\begin{itemize}
% DISS-005-B2234
\item Если подавать модели множество неразмеченных данных, и просить её предсказать будущий токен (слово, точку и т.д.), то она извлекает знание / понимание задачи / мира. Исходно это пошло от моделей автоенкодинга, которым подавалась, например, картинка, далее было множество слоёв, сжимающих картинку, а потом похожая цепочка, которая «разворачивала» картинку. И в итоге её можно обучать прост на исходных данных. А центральный вектор, находящийся по середине, является абстрактным представлением входных данных и может далее использоваться для любой иной задачи. По очень похожему пути обучаются большие языковые модели, ток без полной цепочки воспроизведения. Там на вход модели подаётся весь предыдущий известный текст, а на выходе просят предсказать следующее слово. И за счёт этого – модель обучается представлению ВСЕГО мира, что можем видеть на основе ChatGPT, безусловно, есть множество других хитрых точек, но основная идея такая. Как раз эта идея применима и к нашей задаче, так как в электроэнергетике мы обрабатываем последовательности рядов, и как бы пытаемся предсказать «будущее», чтобы чуть раньше сработать.

% DISS-005-B2235
\item Из выпуска 2.5 (Полная история GPT, блок 2.5: законы масштабирования LLM (Scaling Laws) | Котенков Игорь\commentnote{DISS-005-C0143}{Алексей Евдаков}{\url{https://www.youtube.com/watch?v=J5nGTnv3IN8}}) появились полезные знания о моделях CPT, которые применимы ко всем моделям.:

\begin{itemize}
% DISS-005-B2236
\item Большие модели лучше справляются с задачами, если не оценивать влияние переобучения.
% DISS-005-B2237
\item Больше данных даёт лучшую среднюю точность, если не оценивать влияние излишней зашумлённости
% DISS-005-B2238
\item Лучшие данные, дают лучшую точность (но это старое и очевидное)
% DISS-005-B2239
\item Трансформеры работают лучше рекуррентных сетей для оценки большего контекста. (ещё нужно разобрать, насколько применимо к нам, да и это про сложные модели)
% DISS-005-B2240
\item Статья OpenAI от "kaplan et all"\commentnote{DISS-005-C0144}{Алексей Евдаков}{Найти её} [2020] При увеличении модели в 8 раз, надо в 5 раз больше данных. НО — это не верно
% DISS-005-B2241
\item Статья DeepMind от "Hoffmann et all [20222]" При увеличении модели в 2 раз, надо в 2 раз больше данных (линейно,). Они провели опыт «шиншила», более статистически значимый, там больше моделей и т.д., см картинку ниже (ХХ). Основной их вывод отличающий от OpenAI в том, что можно брать модельки поменьше, но надо больше данных, чем предсказывалось раньше
\end{itemize}
\end{itemize}

% DISS-005-B2242
\begin{figure}[htbp]
\centering
\includegraphics[width=0.9\linewidth]{image97.png}
\end{figure}

% DISS-005-B2243
\par\vspace{0.2cm}

\begin{itemize}
% DISS-005-B2244
\item Может ещё что-то напишу
\end{itemize}

% DISS-005-B2245
\par\vspace{0.2cm}

% DISS-005-B2246
\subsection*{Подготовка датасета}

% DISS-005-B2247
Для более хорошего обучения можно будет использовать и синтетические данные.\par\vspace{0.3cm}

% DISS-005-B2248
Первичную разметку датасета можно сделать на основе дискретных сигналов от нашего РНМ (и/или от других производителей). Поведение такого РНМ не является идеальным, но достаточно хорошее и релевантное. А в будущем можно уже его модифицировать (улучшать разметку в спорных моментах).\par\vspace{0.3cm}

% DISS-005-B2249
\par\vspace{0.2cm}

% DISS-005-B2250
\subsection*{Перенос идей из LLM в текущие задачи}

% DISS-005-B2251
Идея все получил на основе просмотра видео о современных LLM и чтении статей на Хабр. \par\vspace{0.3cm}

% DISS-005-B2252
Я хочу научить автоэнкодировщик (фото 2), это позволит обучать модель на полностью не размеченных данных. По сути, так и делают все большие модели. Как минимум на некоторых шагах. У них нет размеченных данных, но зато очень много данных. И за счёт этого модель обучается понимать, извлекать информацию из самих данных. А потом, на этом – уже создать конкретный классификатор или что-то подобное (что требуется). Причём, модель может сама «создавать данные», за счёт того, что будет сама предсказывать следующий латтент (фото 1,2,3)\par\vspace{0.3cm}

% DISS-005-B2253
Это аналогия того, как учили ChatGPT, взял эту идею вот из этой статьи: \url{https://habr.com/ru/articles/794566/}  (В ней много ссылок и на другие работы, очень хорошая в целом)\par\vspace{0.3cm}

% DISS-005-B2254
Плюс дальше идеи, которые записывал по ходу чтения:\par\vspace{0.3cm}

\begin{itemize}
% DISS-005-B2255
\item я раньше делал пометку на видео в ютубе от этого же человека. (Найти видео). Речь о авторе канала «Сиолошная»
% DISS-005-B2256
\item у него есть и статья о том как это делает ChatGPT
% DISS-005-B2257
\par\url{https://habr.com/ru/companies/ods/articles/716918/}\par\vspace{0.2cm}
% DISS-005-B2258
\item что-то интересное про это может быть в моей в статье по философии.
% DISS-005-B2259
\item может быть полезна точная информация о подстанции. Но... Как это задавать, и в нынешних рамках - нереально. Так что пока, аналог ChatGPT, который смотрит просто текст.
% DISS-005-B2260
\item Увеличение моделей может в разы улучшить результаты, закон масштабирования (Фото 4)
% DISS-005-B2261
\item Обучать модель актуально на данных, где она сильно ошибается, это важный факт дообучения (Фото 5)
% DISS-005-B2262
\item модель предсказывающая будущее и нужна. Вот только вопрос... Что она учится предсказывать следующую точку. А ведь надо уметь делать прогноз и на далëкое будущее, как и человек. Но это сложно, и пока не знаю как. (Хотя, появилось мысль - моделировать на некоторое время вперёд и на этом основании делать выводы - ведь это симулятор мира)
% DISS-005-B2263
\item нейронка - может быть симулятором мира, в некотором смысле аналогом Матлаба, ртдс и прочего.
% DISS-005-B2264
\item можно добавить идею о том, что наш мозг так же учится создавать модель мира, она может быть достаточно проста - например чайка распознаëт маму по Красной линии с белыми точками (клюв). Значит мы так учимся анализировать осциллограммы. (Найти в видео от ALL, он на учёного ссылается, смотрел 02.03.2024)
\end{itemize}

% DISS-005-B2265
\par\vspace{0.2cm}
